\section{数学形式化的历史与生态}

\subsection{从希腊到计算机：形式化的世代}
数学证明始于希腊，欧几里得开创公理化；Hilbert 将形式化写入语言。从符号化到计算时代，数学逐步借助自动化验证：Gentzen 的归纳证明、Curry-Howard 对应、Coq、Isabelle 逼近“人类可读+机器可验证”的交汇。
\begin{itemize}
    \item 500 BCE：希腊几何绘制 axiom-theorem 网络；
    \item 1900s：Hilbert 计划与形式公理体系；
    \item 1970s：Curry-Howard 让类型即证明；
    \item 2010s：Lean 与 Mathlib 拥抱开放协作，目标“人人都能 formalize”；
    \item 2020s：AlphaProof 将 RL 搜索引擎接入 Lean，以应对 IMO 级别命题。
\end{itemize}

Hubert 还强调，数学形式化并不是一条单向的历史，而是“多代 parallel 更新”：Mathlib 的每个 PR 都被 replay 在 RL 环境中，RL agent 反过来又将失败 trace 反馈给 Mathlib，而 governers 在 GitHub/CI 上形成可解释的讨论，这种交叉反馈推动新公理的采纳。

\subsection{Lean、Mathlib 与规范化}
\begin{importantbox}{Lean + Mathlib 的三大原则}
1）\textbf{可复现}：每次定理都有可编译的脚本；2）\textbf{统一}：同一概念只进一个库，避免定义碎片化；3）\textbf{开放}：Mathlib 通过 PR/CI 把贡献者变成治理节点。Hubert 强调：RL agent 既要读懂 Mathlib，也要对 CI 失败敏感。
\end{importantbox}

这些原则直接体现在 AlphaProof 的部署中：formalizer 与 Mathlib 的 PR/CI pipeline 对齐，每次 agent 生成新的 lemma 都要先过 nightly runner；当 CI failed 时，agent 自动降级 search policy，转而将 trace 分享给 human-in-the-loop 组。

Lean 为 AlphaProof 提供的语义元素：
\begin{enumerate}
    \item \textbf{环境 state}：当前定理、前提与 tactic 栈；
    \item \textbf{动作 catalog}：不同 tactic 构成 action space，可定制 BFS、MCTS、Heuristic；
    \item \textbf{反馈}：证明成功就是 +1，failure 会 trigger trace/log；
    \item \textbf{Observation}：Lean 会返回 proof term、goal 树、type 约束。
\end{enumerate}

\subsection{形式化生态中的协作机制}
Hubert 描述了 Lean 社区的协作闭环：Mathlib 维护者通过 PR/build pipeline 控制质量，AlphaProof 团队把“proof search logs + failed tactics”上传到 governance board，形成 RL trace 的输入。好几个 Lean expert 会打开 replay log 以 human-in-the-loop 方式修复 hallucination。

\begin{knowledgebox}{协作式形式化的关键节点}
1）Crowdsourced replay log：记录 agent 失败的 sequence；2）Fairness dashboard：验证 search 是否偏向某类策略；3）Human review loop：当 RL 触发 drift gate 时，expert 可 rollback knowledge base。
\end{knowledgebox}

\subsection{本章小结}
数学形式化的发展、Lean/Mathlib 的规范化原则与 community-driven 协作机制共同支撑 AlphaProof 的上下文，这一节解释了为什么把 RL 连接到 Lean 会带来“可验证的创造力”。
\newpage

